\documentclass[12pt,a4paper]{article}

\usepackage[utf8]{inputenc}
\usepackage[T1]{fontenc}
\usepackage[polish]{babel}
\usepackage{geometry}
\geometry{a4paper, margin=2.5cm}
\usepackage{graphicx}
\usepackage{hyperref}
\usepackage{listings}
\usepackage{xcolor}
\usepackage{float}
\usepackage{amsmath}
\usepackage{amssymb}
\usepackage{booktabs}
\usepackage{tabularx}

\usepackage{titlesec}
\usepackage{tocloft}
\usepackage{changepage} 

\titlespacing*{\subsection}{1.5cm}{3.25ex plus 1ex minus .2ex}{1.5ex plus .2ex}

\setlength{\cftsubsecindent}{1.5cm}

\definecolor{termbg}{rgb}{0.05, 0.05, 0.05}
\definecolor{termfg}{rgb}{0.9, 0.9, 0.9}
\definecolor{codebg}{rgb}{0.95, 0.95, 0.95}
\definecolor{commentgreen}{rgb}{0,0.6,0}
\definecolor{keywordblue}{rgb}{0.13,0.13,1}

\lstdefinestyle{terminal}{
    backgroundcolor=\color{termbg},
    basicstyle=\ttfamily\color{termfg}\footnotesize,
    breaklines=true,
    captionpos=b,
    frame=single,
    rulecolor=\color{gray},
    keywordstyle=\color{white}\bfseries,
    morekeywords={ZALICZONY, Oczekiwano, Otrzymano, TEST},
    literate={ą}{{\k{a}}}1 {ć}{{\'c}}1 {ę}{{\k{e}}}1 {ł}{{\l}}1 {ń}{{\'n}}1 {ó}{{\'o}}1 {ś}{{\'s}}1 {ź}{{\'z}}1 {ż}{{\.z}}1
}

\lstdefinestyle{c-code}{
    language=C,
    backgroundcolor=\color{codebg},
    basicstyle=\ttfamily\footnotesize,
    keywordstyle=\color{keywordblue}\bfseries,
    commentstyle=\color{commentgreen}\itshape,
    stringstyle=\color{red},
    breaklines=true,
    frame=single,
    captionpos=b,
    literate={ą}{{\k{a}}}1 {ć}{{\'c}}1 {ę}{{\k{e}}}1 {ł}{{\l}}1 {ń}{{\'n}}1 {ó}{{\'o}}1 {ś}{{\'s}}1 {ź}{{\'z}}1 {ż}{{\.z}}1
}

\lstdefinestyle{java-code}{
    language=Java,
    backgroundcolor=\color{codebg},
    basicstyle=\ttfamily\footnotesize,
    keywordstyle=\color{keywordblue}\bfseries,
    commentstyle=\color{commentgreen}\itshape,
    stringstyle=\color{red},
    breaklines=true,
    frame=single,
    captionpos=b
}

\begin{document}

\begin{titlepage}
    \begin{center}
        \vspace*{1cm}
        
        {\LARGE \textbf{Politechnika Warszawska}} \\[0.2cm]
        {\Large Wydział Elektryczny} \\[2cm]
        
        {\large Kierunek: Informatyka stosowana} \\[0.1cm]
        {\large Przedmiot: Języki i metody programowania 2} \\[3.5cm]
        
        {\Huge \textbf{Dokumentacja Końcowa Projektu}} \\[1cm]
        {\Large Projekt C: Wyznaczanie współrzędnych węzłów dla wizualizacji\\ grafu planarnego} \\[3.5cm] 
        
        {\large \textbf{Autorzy:}} \\[0.2cm]
        {\large Patryk Pawlak 343367} \\[0.1cm]
        {\large Krzysztof Dębski 343317}
        
        \vfill
        
        {\large Warszawa, 19 kwietnia 2026}
        \vspace*{0.5cm}
    \end{center}
\end{titlepage}

\tableofcontents
\newpage

\section{Wstęp}
Niniejszy dokument stanowi pełną specyfikację końcową modułu obliczeniowego \textbf{GraphVisualizer (C Core)}, realizowanego w ramach przedmiotu Języki i metody programowania 2. 

Celem aplikacji jest automatyczne wyznaczanie współrzędnych $(x, y)$ wierzchołków grafu planarnego w celu jego czytelnej wizualizacji. Aplikacja C działa jako bezstanowy silnik CLI, który przyjmuje topologię grafu, wykonuje intensywne obliczenia rozmieszczające wierzchołki i zwraca gotowe współrzędne. Została zaprojektowana z myślą o pełnej integracji z modułem GUI tworzonym w języku Java.

Dokumentacja ta łączy opisy funkcjonalne, implementacyjne oraz szczegółowe instrukcje pozwalające zespołowi docelowemu na sprawną integrację z silnikiem, bez konieczności analizowania kodu źródłowego w języku C.

\section{Architektura Systemu i Implementacja}

Aplikacja została napisana w standardzie C11 z wykorzystaniem kompilatora GCC. Architektura systemu ściśle przestrzega zasad \textit{Clean Code}, \textit{Minimal Global State} oraz \textit{Defensive Programming}.

\noindent
\begin{minipage}{\textwidth}
\subsection{Struktura Modułowa}
\begin{adjustwidth}{1.5cm}{0cm}
Projekt dzieli się na cztery główne warstwy:
\begin{itemize}
    \item \texttt{core/} – Kontroler sterujący przepływem programu. Przetwarza flagi CLI (przy użyciu biblioteki \texttt{getopt.h}) i agreguje działanie całego programu.
    \item \texttt{io/} – Moduł wejścia/wyjścia. Obejmuje obsługę formatu tekstowego (wczytywanie grafu) oraz eksport do formatów tekstowych i binarnych.
    \item \texttt{algorithms/} – Główny rdzeń obliczeniowy zawierający logikę algorytmów Fruchtermana-Reingolda oraz Tutte'a.
    \item \texttt{utils/} – Zbiór pomocniczych funkcji matematycznych, m.in. implementacja eliminacji Gaussa dla algorytmu barycentrycznego.
\end{itemize}
\end{adjustwidth}
\end{minipage}
\vspace{0.5cm}

\noindent
\begin{minipage}{\textwidth}
\subsection{Kluczowe Struktury Danych}
\begin{adjustwidth}{1.5cm}{0cm}
Stan grafu przechowywany jest w dynamicznie alokowanych strukturach, zarządzanych przez strukturę \texttt{Graph}.
\begin{lstlisting}[style=c-code, caption={Uproszczony model struktur danych}]
typedef struct {
    int id;
    double x, y;
} Node;

typedef struct {
    char name[MAX_NAME_LEN];
    int n_1, n_2;
    double weight;
} Edge;

typedef struct {
    Node *nodes;
    int node_count;
    int node_cap;
    Edge *edges;
    int edge_count;
    int edge_cap;
} Graph;
\end{lstlisting}
\end{adjustwidth}
\end{minipage}
\vspace{0.5cm}

\noindent
\begin{minipage}{\textwidth}
\subsection{Specyfikacja Wewnętrznego API (Wybrane funkcje)}
\begin{adjustwidth}{1.5cm}{0cm}
Interfejs zarządzania grafem (zdefiniowany w \texttt{src\_c/io/graph.h}) oferuje m.in. następujące metody:
\begin{itemize}
    \item \texttt{Graph* create\_graph(void)} – Inicjalizacja pustej struktury. Zwraca wskaźnik lub \texttt{NULL} w przypadku błędu alokacji.
    \item \texttt{int graph\_load\_text(Graph *g, const char *path)} – Wczytuje listę krawędzi. Zwraca kod błędu w przypadku problemów z plikiem/formatem.
    \item \texttt{int graph\_save\_binary(const Graph *g, const char *path)} – Eksportuje wynik działania algorytmów do binarnego pliku wymiany.
    \item \texttt{void free\_graph(Graph *g)} – Zwalnia wszystkie zaalokowane zasoby, zapobiegając wyciekom pamięci.
\end{itemize}
\end{adjustwidth}
\end{minipage}
\vspace{0.5cm}

\section{Instrukcja Użytkownika (CLI)}

Silnik C jest w pełni sterowany z poziomu wiersza poleceń. Pozwala to na łatwe wywoływanie go z poziomu Javy.

\noindent
\begin{minipage}{\textwidth}
\subsection{Argumenty Uruchomieniowe}
\begin{adjustwidth}{1.5cm}{0cm}
Program obsługuje następujące flagi:
\begin{table}[H]
\centering
\begin{tabularx}{\linewidth}{|c|c|X|}
\hline
\textbf{Flaga} & \textbf{Wymagane} & \textbf{Opis} \\ \hline
\texttt{-i <path>} & \textbf{TAK} & Ścieżka do pliku wejściowego (tekstowego) z opisem grafu. \\ \hline
\texttt{-o <path>} & NIE & Plik wyjściowy ze współrzędnymi. Domyślnie wynik wypisywany jest na konsolę (stdout). \\ \hline
\texttt{-a <alg>}  & NIE & Wybór algorytmu rozmieszczania: \texttt{fr} (Fruchterman-Reingold - domyślny) lub \texttt{tutte} (Tutte). \\ \hline
\texttt{-n <int>}  & NIE & Liczba iteracji dla algorytmu FR (domyślnie: 200). Nie dotyczy algorytmu Tutte'a. \\ \hline
\texttt{-b}        & NIE & Wymusza zapis do formatu binarnego. \\ \hline
\texttt{-v}        & NIE & Tryb verbose. Wyświetla diagnostykę i postęp na stderr. \\ \hline
\texttt{-h}        & NIE & Wyświetla ekran pomocy i kończy działanie programu. \\ \hline
\end{tabularx}
\caption{Argumenty CLI programu}
\end{table}

\textbf{Przykłady wywołania:}
\begin{lstlisting}[style=terminal, caption={Przykłady użycia CLI}]
# Standardowe przeliczenie algorytmem FR i wyrzucenie do pliku tekstowego
./graph_layout -i wejscie.txt -o wynik.txt

# Obliczenia algorytmem Tutte'a i eksport binarny dla interfejsu Java
./graph_layout -i wejscie.txt -o wynik.bin -a tutte -b

# Uruchomienie z podgladem postepu (verbose)
./graph_layout -i wejscie.txt -a fr -n 500 -v
\end{lstlisting}
\end{adjustwidth}
\end{minipage}
\vspace{0.5cm}

\section{Format Danych}

\noindent
\begin{minipage}{\textwidth}
\subsection{Plik Wejściowy (Tekstowy)}
\begin{adjustwidth}{1.5cm}{0cm}
Format wejściowy stanowi płaską listę krawędzi oddzielonych znakami białymi. Każda krawędź posiada nazwę, identyfikatory węzłów, które łączy oraz wagę.

Format jednej linii: \texttt{<nazwa\_krawedzi> <wierzcholek\_A> <wierzcholek\_B> <waga\_krawedzi>}

\textbf{Przykład pliku \texttt{input.txt}:}
\begin{lstlisting}[style=terminal]
E1 1 2 1.0
E2 2 3 1.5
E3 3 1 1.0
E4 3 4 2.0
\end{lstlisting}
\end{adjustwidth}
\end{minipage}
\vspace{0.5cm}

\noindent
\begin{minipage}{\textwidth}
\subsection{Plik Wyjściowy (Tekstowy)}
\begin{adjustwidth}{1.5cm}{0cm}
Jeśli nie użyto flagi \texttt{-b}, program wygeneruje plik tekstowy. Zawiera on wyliczone pozycje $(x,y)$ unikalnych wierzchołków z dokładnością do 6 cyfr po przecinku.

Format jednej linii: \texttt{<wierzcholek> <x> <y>}

\textbf{Przykład pliku \texttt{output.txt} (dla powyzszego wejścia):}
\begin{lstlisting}[style=terminal]
1 0.000000 0.000000
2 1.000000 0.000000
3 0.500000 0.866000
4 0.500000 1.866767
\end{lstlisting}
\end{adjustwidth}
\end{minipage}
\vspace{0.5cm}

\section{Specyfikacja Mostka Binarnego (C-Java)}
W celu zapewnienia wysoce wydajnej wymiany danych między rdzeniem obliczeniowym C a interfejsem graficznym Java, głównym formatem produkcyjnym (flaga \texttt{-b}) jest zoptymalizowany plik binarny. 


\noindent
\begin{minipage}{\textwidth}
\subsection{Struktura Rekordu Binarnego}
\begin{adjustwidth}{1.5cm}{0cm}
Każdy węzeł grafu zapisywany jest sekwencyjnie. Nie występują tu żadne nagłówki plików. Rozmiar pojedynczego rekordu wynosi dokładnie \textbf{20 bajtów}:
\begin{itemize}
    \item \texttt{int id} (4 bajty) – Unikalny identyfikator węzła.
    \item \texttt{double x} (8 bajtów) – Współrzędna X zapisana w standardzie IEEE 754.
    \item \texttt{double y} (8 bajtów) – Współrzędna Y zapisana w standardzie IEEE 754.
\end{itemize}

Spójność danych: Całkowity rozmiar pliku wyjściowego zawsze będzie równy $N \times 20$ bajtów, gdzie $N$ to liczba unikalnych wierzchołków w przetworzonym grafie.
\end{adjustwidth}
\end{minipage}
\vspace{0.5cm}

\noindent

\vspace{0.5cm}

\section{Zaimplementowane Algorytmy Rozmieszczania}

System udostępnia dwa niezależne mechanizmy wyznaczania współrzędnych. Aplikacja Java powinna umożliwić użytkownikowi przełączanie między nimi (flaga \texttt{-a}).

\noindent
\begin{minipage}{\textwidth}
\subsection{Algorytm Fruchtermana-Reingolda (Podejście siłowe, \texttt{-a fr})}
\begin{adjustwidth}{1.5cm}{0cm}
Uniwersalny algorytm traktujący graf jak system fizyczny. 
\begin{itemize}
    \item \textbf{Zasada działania:} Krawędzie grafu traktowane są jak sprężyny przyciągające wierzchołki, natomiast same wierzchołki działają na siebie jak naładowane elektrycznie cząsteczki (siły odpychania).
    \item \textbf{Cechy szczególne:} Obliczenia wykorzystują wagi krawędzi (wpływają na siłę przyciągania). Posiada mechanizm chłodzenia zapewniający stabilizację w kolejnych iteracjach. System wymusza ramkę, zapobiegając ucieczce węzłów w nieskończoność.
    \item \textbf{Zastosowanie:} Radzi sobie z grafami niespójnymi oraz nieplanarnymi. Idealny dla większości ogólnych przypadków użycia.
\end{itemize}
\end{adjustwidth}
\end{minipage}
\vspace{0.5cm}

\noindent
\begin{minipage}{\textwidth}
\subsection{Algorytm Tutte'a (Twierdzenie o sprężynach, \texttt{-a tutte})}
\begin{adjustwidth}{1.5cm}{0cm}
Algorytm barycentryczny dedykowany grafom planarnym, gwarantujący brak przecięć krawędzi.
\begin{itemize}
    \item \textbf{Zasada działania:} Znajduje cykl w grafie, a następnie "przypina" go do obwodu geometrycznego wielokąta wypukłego. Pozycje pozostałych (wewnętrznych) wierzchołków wyznaczane są poprzez rozwiązywanie układu równań liniowych (metoda Gaussa), tak aby każdy wierzchołek znalazł się w środku ciężkości swoich sąsiadów.
    \item \textbf{Mechanizmy obronne:} Program automatycznie używa algorytmu DFS do znalezienia zewnętrznego cyklu. Dodatkowo wykorzystuje BFS do walidacji spójności grafu – \textbf{warunek konieczny}.
    \item \textbf{Zastosowanie:} Służy wyłącznie do grafów spójnych i planarnych (oraz posiadających cykl). W przypadku podania grafu niespójnego program zwróci kod błędu.
\end{itemize}
\end{adjustwidth}
\end{minipage}
\vspace{0.5cm}

\section{Sytuacje Wyjątkowe i Obsługa Błędów}

Program stosuje rygorystyczne podejście defensywne. Zamiast powodować awarię (segmentation fault), w sytuacjach brzegowych zwalnia zasoby i zwraca odpowiedni kod zakończenia.

\begin{table}[H]
\centering
\begin{tabularx}{\textwidth}{|c|l|X|}
\hline
\textbf{Kod} & \textbf{Typ Błędu} & \textbf{Przykłady / Opis} \\ \hline
\texttt{0} & SUKCES & Wykonanie poprawne. Wyniki zapisane do pliku. \\ \hline
\texttt{1} & BŁĄD PLIKU (\texttt{ERR\_FILE}) & Plik wejściowy nie istnieje lub brak praw dostępu. Brak uprawnień do zapisu pliku wyjściowego. \\ \hline
\texttt{2} & BŁĄD FORMATU (\texttt{ERR\_FORMAT})& Niezgodna liczba kolumn w pliku wejściowym (np. 2 zamiast 4). Znaki niebędące liczbami tam, gdzie oczekiwano double/int. \\ \hline
\texttt{3} & BŁĄD ALGORYTMU (\texttt{ERR\_ALGO})& Uruchomienie algorytmu Tutte'a dla grafu bez cyklu lub dla grafu niespójnego. Prowadziłoby to do dzielenia przez zero w macierzy. \\ \hline
\texttt{4} & BŁĄD ARGUMENTÓW (\texttt{ERR\_ARGS})& Brak flagi \texttt{-i}. Podanie nieznanego parametru flagi (np. \texttt{-a xyz}). \\ \hline
\end{tabularx}
\caption{Specyfikacja kodów powrotu silnika}
\end{table}

\section{System Testowy i Walidacja Jakości}

Aby zapewnić stabilność i przewidywalność pracy programu C, zaimplementowano zautomatyzowane ramy testowe. Dodatkowo użyto narzędzia \texttt{Valgrind} do kontroli pamięci.

\noindent
\begin{minipage}{\textwidth}
\subsection{Środowisko i Metodologia}
\begin{adjustwidth}{1.5cm}{0cm}
Wykorzystano skrypt \texttt{generator.c}, który masowo produkuje przypadki brzegowe i obciążeniowe. Test runner analizuje kody wyjścia i weryfikuje ich zgodność z założeniami. Przeprowadzono m.in.:
\begin{itemize}
    \item Testy stabilności FR: Obciążenia grafami $>500$ wierzchołków dla 1000 iteracji.
    \item Testy spójności Tutte'a: Wstrzykiwanie grafów losowych w celu sprawdzenia reakcji modułu weryfikacyjnego.
    \item Weryfikacja precyzji: Ekstremalne wielkości wejściowe (wagi rzędu $10^{-6}$ do $10^7$).
\end{itemize}
\end{adjustwidth}
\end{minipage}
\vspace{0.5cm}

\noindent
\begin{minipage}{\textwidth}
\subsection{Raport z Wykonania Zestawu Testowego}
\begin{adjustwidth}{1.5cm}{0cm}
Poniżej znajduje się udokumentowany zrzut logu ze zautomatyzowanego runnera. Potwierdza on 100\% skuteczność obsługi sytuacji wyjątkowych.

\begin{lstlisting}[style=terminal, caption={Log z konsoli z weryfikacji automatycznej}]
========================================
 ROZPOCZYNAM AUTOMATYCZNE TESTY SYSTEMU
========================================

--- TEST: Obciazeniowy - 100 wezlow (Tutte) ---
[ZALICZONY] Oczekiwano: 0, Otrzymano: 0

--- TEST: Obciazeniowy - 500 wezlow (FR) ---
[ZALICZONY] Oczekiwano: 0, Otrzymano: 0

--- TEST: Krawedziowy - Zbyt maly graf (Tutte musi przerwac - Kod 3) ---
[ZALICZONY] Oczekiwano: 3, Otrzymano: 3

--- TEST: Krawedziowy - Ekstremalne wagi (Precyzja double) ---
[ZALICZONY] Oczekiwano: 0, Otrzymano: 0

--- TEST: Weryfikacja Valgrind (Brak wyciekow pamieci) ---
[ZALICZONY] Oczekiwano: 0, Otrzymano: 0

========================================
WYNIK TESTOW: 5 / 5 ZALICZONYCH
WSZYSTKIE TESTY ZAKONCZONE SUKCESEM!
========================================
\end{lstlisting}
\end{adjustwidth}
\end{minipage}
\vspace{0.5cm}

\noindent
\begin{minipage}{\textwidth}
\subsection{Gwarancja braku wycieków (Valgrind)}
\begin{adjustwidth}{1.5cm}{0cm}
Jak prezentuje log, w każdym przypadku poprawnym oraz błędnym program wywoływał procedurę \texttt{free\_graph()} oraz odpowiednie sprzątanie zasobów plikowych \texttt{fclose()}. Raport pamięci Valgrind ("Clean memory report") potwierdza \textbf{zero utraconych bajtów}, co oznacza, że wielokrotne wywoływanie aplikacji w tle (proces po procesie) nie zdestabilizuje środowiska operacyjnego użytkownika.
\end{adjustwidth}
\end{minipage}
\vspace{0.5cm}

\section{Wnioski}
Silnik obliczeniowy \textbf{GraphVisualizer (C Core)} osiągnął status ukończony (wersja produkcyjna). Spełnia wszystkie założenia specyfikacji funkcjonalnej i implementacyjnej.

Zaimplementowane algorytmy poprawnie rozmieszczają węzły grafów. Architektura programu oparta na bezstanowym CLI i formacie binarnym udostępnia wysoce wydajne, klarowne API pozwalające na bezproblemową budowę docelowego GUI w środowisku Java. Zabezpieczenia matematyczne zapobiegające awariom na niepoprawnych danych wejściowych czynią program odpornym na błędy, co udowodniono w udokumentowanych, bezbłędnych przebiegach testów automatycznych.

\end{document}